\section{Preliminaries and problem formulation}
\label{sec:prelim}

We model hybrid RAG pipelines as plans over a small operator algebra,
state what a quality contract obliges a service to do, and then
distinguish two ways of discharging that obligation -- by inference,
which is what the literature does, and by accounting, which is what we
propose.

\subsection{Hybrid RAG plans over heterogeneous knowledge sources}
\label{sec:algebra}

Let $\mathcal{S}$ be a set of evidence sources: relational tables, text
corpora with sparse and dense indexes, and remote structured knowledge
APIs. A \emph{plan} $p$ is a directed acyclic graph built from typed
operators
\[
\textsf{Retrieve}_{s,k} \mid \textsf{Fuse} \mid \textsf{Rerank}_m
\mid \textsf{Filter}_\phi \mid \textsf{Generate}_g ,
\]
where $\textsf{Retrieve}_{s,k}$ fetches $k$ evidence units from source
$s$ (SQL row selection, BM25, dense approximate nearest neighbour, or
API calls), $\textsf{Fuse}$ merges rankings, $\textsf{Rerank}_m$
re-scores with a cross-encoder $m$, $\textsf{Filter}_\phi$ drops
evidence failing a predicate $\phi$, and $\textsf{Generate}_g$ produces
the answer with generator $g$. Executing $p$ on a query $q$ yields an
answer, an evidence set with a citation map, and a realised cost
metered at run time.

The plan space is the Cartesian product of an \emph{access path} -- the
retrieval-side configuration, which is shared across generators once
materialised -- and a \emph{generator}, which may be a hosted commercial
model or a locally served open-weight one. Systems in the literature
almost universally restrict attention to a diagonal of this product, an
escalation ladder in which retrieval depth and generator tier increase
together. Section~\ref{sec:results} shows the diagonal misses the plans
a strict contract needs, and also that enlarging the space is not by
itself an improvement, because it aggravates the multiplicity problem
that the incumbent certification machinery has.

\subsection{Evidence contracts}

\begin{definition}[Evidence contract]
\label{def:contract}
An \emph{evidence contract} is a pair $\langle \ell, \alpha\rangle$
where $\ell : (a, q) \mapsto [0,1]$ is a bounded, auditable violation
loss and $\alpha \in (0,1)$ is the admissible violation rate. Examples
of $\ell$ are
$\mathbf{1}[\mathrm{quality}(\cdot) < \tau_q]$,
$\mathbf{1}[\mathrm{citation\ recall}(\cdot) < \tau_c]$,
$\mathbf{1}[\max_{e} \mathrm{age}(e) > \Delta_f]$, and
$\mathbf{1}[\mathrm{latency}(\cdot) > B_\ell]$.
\end{definition}

A financial-services assistant is the running example: at most $10\%$ of
answers incorrect, at most $10\%$ carrying an unsupported claim, at
minimal spend. The requirement is on the \emph{service}, over a stream
of queries, not on any single answer -- which is what makes it a rate
and what will let us enforce it by bookkeeping.

\subsection{The action space, and the action the literature omits}
\label{sec:actions}

The executor's action set is
$\mathcal{A} = \mathcal{P} \cup \{\bot\}$, where $\mathcal{P}$ is the
plan space and $\bot$ is the decision to \emph{decline}: return
"insufficient evidence", route the query to a human, or degrade to a
document list. Declining produces no answer and therefore no incorrect
answer, so
\begin{equation}
\label{eq:abstain}
\ell(\bot, q) = 0 \quad\text{for every } q ,
\qquad c(\bot, q) = c_\bot ,
\end{equation}
where $c_\bot$ prices the handoff. We treat $c_\bot$ as a parameter of
the deployment rather than fixing it, and sweep it, because its value is
an operational fact about the service and not something a paper can
assert: a consumer search box degrades cheaply, a clinical decision
support system does not.

Omitting $\bot$ is the single modelling decision that makes strict
contracts unreachable. With $\mathcal{P}$ alone the attainable risk is
bounded below by $\min_p r_p$, so a contract stricter than the best
available pipeline is infeasible by construction rather than by
statistics -- and $\min_p r_p$ is $0.59$ on CRAG and QAMPARI, $0.27$ on
HybridQA and $0.12$ on ASQA, so the $10\%$ contract of our running
example is unexpressible on all four.
Abstention is standard in selective prediction
\citep{geifman2017selective,angelopoulos2021gentle} and in conformal
methods, but it is treated there as a property of a single predictor;
here it is an \emph{action} the optimizer costs and schedules alongside
the plans, which is what makes the risk--cost frontier continuous.

\subsection{Two ways to discharge the obligation}

Let $R(\pi) = \E_{q}[\ell(\pi,q)]$ and
$\mathrm{cost}(\pi) = \E_{q}[c(\pi,q)]$. The design goal is
\begin{equation}
\label{eq:problem}
\min_{\pi} \ \mathrm{cost}(\pi)
\quad \text{s.t.} \quad R(\pi) \le \alpha ,
\end{equation}
with $\mathcal{D}$ unknown. Since $R$ cannot be computed, the
obligation must be discharged some other way, and there are exactly two.

\begin{definition}[Certified deployment]
\label{def:certified}
A procedure mapping a calibration sample $Z_{1:n}$ to a policy
$\hat\pi$ is \emph{$(\alpha,\delta)$-certified} if
$\Prob_{Z_{1:n}}( R(\hat\pi) \le \alpha ) \ge 1-\delta$,
where the probability is over the calibration draw and the guarantee is
about the population $\mathcal{D}$ that generated it.
\end{definition}

\begin{definition}[Contract-enforcing execution]
\label{def:enforcing}
An executor is \emph{contract-enforcing} if the actions
$a_1, a_2, \dots$ it emits on a query stream
$q_1, q_2, \dots$ satisfy
\[
\textstyle\sum_{s \le t} \ell(a_s, q_s) \;\le\; \alpha t ,
\quad \text{all } t \ge 1 ,
\]
for every query sequence, with no probability attached.
\end{definition}

The two are not stronger and weaker versions of one statement; they
differ in what they are about. Definition~\ref{def:certified} concerns
an unobservable population mean and is conditional on the deployment
traffic being exchangeable with the calibration sample.
Definition~\ref{def:enforcing} concerns the realised trajectory --
the thing an auditor would actually measure -- and is conditional on
nothing. It is also the property an operator asked to sign a
service-level agreement needs, because an SLA is checked on the log,
not on a hypothetical population.

Existing RAG optimizers provide neither. Workload-level optimizers such
as Abacus \citep{russo2025abacus} compare point estimates against the
constraint; adaptive routers \citep{jeong2024adaptiverag} make no
feasibility statement. Conformal RAG methods
\citep{li2024traq,kang2024crag} achieve Definition~\ref{def:certified}
for a single fixed pipeline, leaving the cost-minimisation dimension
of~\eqref{eq:problem} untouched, and Learn-then-Test
\citep{angelopoulos2021ltt} achieves it for a searched policy at a
multiplicity price. This paper achieves
Definition~\ref{def:enforcing}, and recovers the cost objective as a
separate, purely statistical problem in which errors are paid for in
money.

\begin{rmk}[What the guarantee is about]
A contract is stated with respect to a loss the operator can audit.
Both definitions apply to whatever $\ell$ is declared -- token-F1
against gold answers, NLI-based citation entailment, an LLM judge
following a benchmark's official protocol, evidence age, wall-clock
latency. Neither establishes that the declared predicate captures the
application's intent, and neither promises per-query correctness. The
deterministic guarantee additionally requires the loss to be observed,
which costs money; Section~\ref{sec:labels} prices observing only part
of it, and Section~\ref{sec:discussion} discusses a systematically
biased auditor, which no accounting scheme can repair.
\end{rmk}
